\runninghead{Clay, Vine, and Fidelity}
\properlane{covenant-fidelity-and-human-plasticity}{The Potter's Clay, the Vine, and the Faithfulness of God}

Where the first interpretation begins with God's apparent silence, this second interpretation begins with human breakdown. Isaiah confesses that ``all our justices'' have become like a polluted cloth, that Israel has withered like a fallen leaf swept away by iniquity, and that no one even rises up to call upon God's name (Isa.~64:6--7~Vg.). The Responsorial Psalm pictures a vine in ruins that cannot restore itself unless God turns his people back (Ps.~80[79]:4, 15). In the Second Reading, Paul writes to a church about to be rebuked for factions and worldliness (1~Cor.~1:10--12; 3:1--3), yet promises that ``God is faithful'' and will confirm them ``unto the end without crime'' (1~Cor.~1:8--9). If human goodness is so easily stained and human loyalty so fragile, on what ground can the Church hope to stand blameless on the Day of Christ --- self-improvement, or God's covenant faithfulness acting through the Son of man without overriding human freedom? St.~Thomas Aquinas on Isaiah~63--64, St.~Augustine on divine hardening and Psalms~25(24), 80(79), and~85(84), St.~Robert Bellarmine on Psalms~80(79) and~85(84), and St.~John Chrysostom on 1~Corinthians~1:4--9 each press that balance in the passages before us.

\subsection{Stained Righteousness, Clay, and the Vine}

At the opening of the Mass, the psalmist of the Entrance Antiphon \emph{Ad te levavi} (Ps.~25[24]:1--3) turns from self-trust to God alone; St.~Augustine (\work{Enarrationes in Psalmos}, Ps.~24 §§1--2) reads that plea against Adam's pride, which trusted in self and fell into humiliation. With the \work{Gloria} omitted (GIRM~53), the Collect (\emph{Da, quaesumus, omnipotens Deus}; modern text withheld) opens Advent prayer, and its ancient Gelasian antecedent (II.lxxxiv) asks almighty God to bestow the very \emph{voluntatem} (``will'') to go out and meet Christ in righteous deeds.

On Isaiah~63--64, St.~Thomas Aquinas (\work{Expositio super Isaiam ad litteram}, ch.~63--64) takes a different route from St.~Jerome on two key verses. First, where Jerome reads ``made us to err'' and ``hardened our heart'' (63:17) through God's excessive forbearance, Aquinas explains them through permissive withdrawal of aid --- permitting sin, not infusing grace, and not correcting --- citing Romans~9:18. St.~Augustine (\work{De gratia et libero arbitrio}, ch.~45 [XXIII]) holds the same boundary from both sides: God's just judgment in withholding grace never turns God into the author of evil and never cancels human free will, for Pharaoh hardened his own heart even as God hardened it by just sentence. Second, on Isaiah~64:3b~Vg.\ (Heb.~64:2b; \emph{Descendisti, et a facie tua montes defluxerunt}, ``thou didst come down, and at thy presence the mountains melted away''), after linking the omitted clause 64:3a~Vg.\ (Heb.~64:2a) to Israel's dread at Sinai (Exod.~20), Aquinas prefers a historical literal sense over Jerome's immediate Incarnational reading: Israel appeals to God's proven covenant rescue from Egypt (Exod.~4; Ps.~18[17]:10) and the melting of the cliffs of Arnon (Num.~21), treating past mercy as ground for renewed deliverance. Then on 64:6~Vg.\ (Heb.~64:5; \emph{quasi pannus menstruatae universae justitiae nostrae}, ``all our justices as the rag of a menstruous woman''), Aquinas protects authentic graced virtue by naming four concrete stains that can soil a righteous deed: mortal sin, negligence, twisted intention, and pride. Left to itself, clay cannot shape itself; Israel's hope rests on confessing ``we are clay: and thou art our maker'' (Isa.~64:8~Vg.\ [Heb.~64:7]; cf.\ Jer.~18).

In the Responsorial Psalm (Ps.~80[79]), the imagery shifts from clay to a ravaged vineyard, where St.~Augustine (\work{Enarrationes in Psalmos}, Ps.~79 §§9, 11) and St.~Robert Bellarmine (\work{Explanatio in Psalmos}, Ps.~79, explanations~14--18) trace God's faithfulness to his covenant people. God plants no other vine to replace Israel; he visits and perfects his vineyard ``upon the Son of man'' (vv.~15--18). Bellarmine argues that the passage refers literally to Christ --- formed by God's right hand in the Virgin through the Holy Spirit (Luke~1:35) --- because after the Babylonian captivity the Jewish people never recovered their former position, while Augustine identifies the Son of man with Christ Jesus, the sole foundation of the Church (1~Cor.~3:11). Because God must convert his people first (\emph{converte nos}, v.~4), his quickening grace restores the voice that Isaiah~64:7~Vg.\ (Heb.~64:6) lamented had died out, and the psalm closes with the vow ``we will call upon thy name'' (Ps.~80[79]:19).

\subsection{Faithful God, Willing Cooperation, and Graced Work}

In the Second Reading (1~Cor.~1:3--9), St.~John Chrysostom (\work{Homilies on First Corinthians}, Homily~II, §§1--8) examines how Paul applies that theology of grace to a boastful, divided congregation. Paul begins with what seem to be praises, yet as Chrysostom observes, he ties every endowment to \emph{the grace of God} (1:4) and repeats the Name of Jesus Christ relentlessly to break party pride before rebuking their factions in 1:10. Even promising that God will confirm them ``unto the end without crime'' (1:8) quietly tells them they are wavering. On \emph{Fidelis Deus} (``God is faithful,'' 1:9), Chrysostom rejects fatalism: God keeps his promise to bring believers into the fellowship (\emph{societatem}, Greek \emph{koinonia}) of his Son ``if we be not quite impatient of His rein,'' recalling 1~Corinthians~2:9. Cornelius a Lapide (\work{In 1~Cor.}~1:8--9, Vol.~7) states the same equilibrium through the Council of Trent (Session~VI, ch.~11 and~13, citing Augustine): God confirms believers ``so far as His part is'' and will not desert them unless he is first deserted.

Before the Gospel, the assembly sings \emph{Ostende nobis, Domine} (Ps.~85[84]:8), asking for God's mercy rather than pleading human merit --- leaving no room for pride, as Augustine remarks (\work{Enarr.}, Ps.~84 §6). Taken alongside Isaiah's confession of stained deeds and Paul's promise of faithful grace, the authority the householder gives his servants in Mark~13:34 (\emph{potestatem cujusque operis}) is not an unsupported command: St.~Bede (\work{In Marci Evangelium Expositio} IV.13) identifies it explicitly as the grace of the Holy Spirit enabling every good work. Watching and praying therefore mean remaining pliable under the Potter's hand and alert in the Spirit's gift.

Yet Mark~13:33--37 also marks the chief limit of this interpretation. Unlike Isaiah~64 or 1~Corinthians~1, Jesus' parting parable says little directly about sin, clay, or covenant restoration; unless held beside the Gospel's abrupt charge to stay awake, Pauline and Isaian grace theology can mute Mark's urgent call to watchfulness.

\subsection{Gifts Drawn from God's Own Benefits}

In the Prayer over the Offerings (\emph{Suscipe, quaesumus, Domine}; modern text withheld), the ancient Veronense antecedent (XXVIIII.372) states the priority of grace directly: the Church offers back to God gifts drawn from his own bestowed benefits (\emph{quae de tuis offerimus collata beneficiis}), asking that temporal devotion may yield eternal redemption. At the Seasonal Preface of Advent (GIRM~364; modern text withheld), thanksgiving turns to the faithful God whose covenant ways Augustine (\work{Enarr.}, Ps.~24 §10) describes as ``mercy and truth'' across Christ's two Advents. At Holy Communion, the assembly sings \emph{Dominus dabit benignitatem, et terra nostra dabit fructum suum} (Ps.~85[84]:13). Reading that verse at the altar brings Bellarmine's and Augustine's expositions together: Christ, the fruit brought forth from the untouched earth of the Virgin Mary through the Spirit (Bellarmine, Ps.~84, explanation~12), enters hearts furrowed by confession and pours out interior sweetness so that human soil delights in righteousness and bears fruit (Augustine, Ps.~84 §§11--12). The Prayer after Communion (\emph{Prosint nobis}; modern text withheld) closes the Communion Rite, and the Concluding Blessing and dismissal (GIRM~90) send the assembly back to daily life.

\subsection{The Four Senses of Scripture}

\begin{fourSenses}
\item[Literal.] Israel confesses covenant failure, remembers God's deliverance from Egypt and the cliffs of Arnon (Aquinas on Isa.~64:3b~Vg.\ [Heb.~64:2b]), and appeals to its Father and Potter (Isa.~63--64); the psalmist asks the Shepherd to visit and restore his ruined vineyard (Ps.~80[79]); Paul grounds a gifted yet quarrelsome Corinthian church in God's grace and faithfulness (1~Cor.~1:3--9); Jesus entrusts his household and tasks to his servants (Mark~13:33--37).
\item[Allegorical.] The vine perfected upon the Son of man (Ps.~80[79]:16--18, Augustine and Bellarmine) and the earth yielding fruit through divine goodness (Ps.~85[84]:13, Bellarmine) signify the Church built upon Christ born of the Virgin Mary and endowed with the Spirit's gifts (1~Cor.~1:5--7; Mark~13:34, Bede).
\item[Moral.] Because human righteous acts are easily tainted by mortal sin, negligence, perverse intention, or pride (Aquinas on Isa.~64:6~Vg.\ [Heb.~64:5]) and \emph{Fidelis Deus} never forces an unwilling heart (Chrysostom and Lapide on 1~Cor.~1:9), believers practice humble confession, yield like clay to the Potter, and cooperate with the grace that makes righteousness sweet (Augustine on Ps.~85[84]:13).
\item[Anagogical.] The faithful God who calls his people into the fellowship of his Son (1~Cor.~1:9) will confirm those who persevere ``without crime'' on the Day of Christ's coming (1~Cor.~1:8), transforming temporal devotion into eternal redemption (\work{Veronense} antecedent).
\end{fourSenses}
